\section{The actual numerical quotient carries its canonical curve-weak topology}
\label{sec:native-numerical-curve-topology}

\begin{definition}[The whole original curve-degree embedding]
\label{def:native-numerical-curve-weak-topology}
\lean{EnvelopeProperty.NativeRealNumericalClasses.numericalDegreeVector,EnvelopeProperty.NativeRealNumericalClasses.curveWeakTopology}
\leanok
\uses{thm:native-real-numerical-classes}
For an actual scheme projective over its original field, give its actual
real numerical quotient the topology induced by its whole linear map into
$\prod_C\mathbf R$, where $C$ ranges over all actual closed integral curves
and the coordinate is the true full Euler degree on the original curve.
This is the curve-weak topology on that same original quotient.
\end{definition}

\begin{theorem}[Actual separating degrees give a Hausdorff real vector topology]
\label{thm:native-numerical-curve-weak-vector-topology}
\lean{EnvelopeProperty.NativeRealNumericalClasses.numericalDegreeVector_injective,EnvelopeProperty.NativeRealNumericalClasses.numericalDegreeVector_isEmbedding,EnvelopeProperty.NativeRealNumericalClasses.curveWeakT2Space,EnvelopeProperty.NativeRealNumericalClasses.curveWeakTopologicalAddGroup,EnvelopeProperty.NativeRealNumericalClasses.curveWeakContinuousSMul}
\leanok
\uses{def:native-numerical-curve-weak-topology,thm:native-real-numerical-classes}
The whole original degree map is a topological embedding. In particular,
this curve-weak topology is Hausdorff, addition and negation are continuous,
and scalar multiplication by real numbers is jointly continuous.
\end{theorem}
\begin{proof}
All actual curve degrees jointly separate the original numerical quotient,
so the whole degree map is injective. The topology was induced by that
same map, hence it is a topological embedding into the Hausdorff product.
Addition, negation and real scalar multiplication are coordinatewise
continuous on the product. The map preserves these operations as an actual
linear map. Its induced topology therefore makes all these same operations
continuous on the original numerical quotient.
\end{proof}

\begin{theorem}[Actual numerical convergence is precisely convergence of every original curve degree]
\label{thm:native-numerical-curve-weak-convergence}
\lean{EnvelopeProperty.NativeRealNumericalClasses.continuous_iff_curveDegree,EnvelopeProperty.NativeRealNumericalClasses.tendsto_iff_curveDegree,EnvelopeProperty.NativeRealNumericalClasses.continuous_curveDegree}
\leanok
\uses{thm:native-numerical-curve-weak-vector-topology}
A map from any topological space into the original numerical quotient is
continuous if and only if each of its true original curve degrees is
continuous. For any filter, a net converges to an original numerical class
if and only if every actual original curve degree converges to its degree
on that class. In particular, every true original curve-degree functional
is continuous.
\end{theorem}
\begin{proof}
The induced topology tests continuity and convergence after composition
with the whole degree embedding. The product topology tests either condition
at every coordinate. These coordinates are literally the actual original
curve-degree functionals. Applying the continuity assertion to the identity
map proves each functional continuous.
\end{proof}

\begin{theorem}[Every actual finite line-family map is continuous]
\label{thm:native-numerical-line-family-continuous}
\lean{EnvelopeProperty.NativeRealNumericalClasses.continuous_lineFamilyMap,EnvelopeProperty.NativeRealNumericalClasses.tendsto_lineFamilyMap}
\leanok
\uses{thm:native-numerical-curve-weak-convergence,thm:native-numerical-line-family-maps}
For every actual finite family of original invertible lines, the entire
original real coefficient map into the same numerical quotient is
continuous. Convergence of its actual coefficient vectors consequently
gives convergence of their actual original numerical classes.
\end{theorem}
\begin{proof}
On every original curve its degree is the finite sum of each original
coefficient times the fixed true full Euler degree of that same line.
Each coefficient projection is continuous, as is multiplication by this
fixed real degree. Their finite sum is continuous. The preceding
coordinatewise criterion proves continuity of the entire original family
map, and continuity implies the asserted filter convergence.
\end{proof}

\begin{theorem}[The genuine original nef cone is curve-weakly closed]
\label{thm:native-numerical-nef-curve-weak-closed}
\lean{EnvelopeProperty.NativeRealNumericalClasses.isClosed_nefCone}
\leanok
\uses{thm:native-numerical-curve-weak-convergence,thm:native-numerical-nef-cone}
The actual cone cut out by the true degrees on every original closed
integral curve is closed in this same original curve-weak numerical space.
\end{theorem}
\begin{proof}
For each original curve, its actual degree is continuous, so its inverse
image of the closed real half-line $[0,\infty)$ is closed. The genuine
original nef cone is their entire intersection, which is closed.
\end{proof}

\begin{remark}
This construction makes no finite-dimensionality assumption. The numerical
space in the FHX and BJ definitions has its finite-dimensional real vector
space topology. Genuine projective numerical finite dimensionality and its
identification with this canonical curve-weak topology remain mathematical
obligations. Convergence proved here is explicitly curve-weak until that
identification is established; it does not by itself close the full general
numerical-class psh comparison. Varying original-line approximants can now
be defined on this actual topological quotient without restricting them to
a fixed line ray.
\end{remark}
